% ==================== Chapter 3: Requirements and Architecture ====================
\newpage\section{System Requirements and Overall Design} [Estimated: 4,500 words]

\subsection{System Boundaries} [Estimated: 800 words]
Our framework is positioned as a ``non-intrusive, low-power privacy compliance supervision tool for end users''. The boundary constraints are:

\begin{itemize}
\item \textbf{Monitoring only, no interception}: The framework reads historical data without modifying or intercepting permission calls, consistent with Privacy by Design——enhancing transparency without affecting normal system operation.
\item \textbf{No root required}: all data reading uses public APIs (AppOpsManager), ensuring ordinary users can install and use it. AppOpsManager is Android's system service for application operation tracking, serving two primary purposes: access control and tracking. It tracks many important events, including all accesses to runtime permission protected APIs. Importantly, the tracked data can only be read by system components, and our framework accesses this data through the public API surface without requiring system-level privileges.
\item \textbf{Low resource footprint}: WorkManager's background scheduling naturally adapts to system power-saving policies (Doze mode). Target memory <150 MB, extra battery drain <2\%. PeriodicWorkRequest in WorkManager is subject to OS battery optimizations such as Doze mode, and execution may be delayed as the system optimizes resource usage. The minimum interval for periodic work is 15 minutes, which aligns with our design goal of low-frequency auditing.
\item \textbf{Offline operation}: all auditing and decision logic runs locally on the device; no data is transmitted to external servers, preserving user privacy.
\end{itemize}

Functional requirements include: (1) periodic permission call collection, (2) category-based baseline loading, (3) dynamic threshold computation and violation detection, (4) notification generation with tiered severity levels, (5) audit trail persistence and retrieval, and (6) a test harness for automated validation.

The design of the data persistence layer employs Room, Android's recommended abstraction layer over SQLite. Room provides compile-time verification of SQL queries and offers a more concise API compared to direct SQLite usage[reference:7]. Database operations are performed on background threads using Kotlin coroutines to avoid blocking the main thread[reference:8]. The ViolationRecord table stores packageName, permissionType, violationCount, and riskLevel for each detected incident, enabling historical audit trail retrieval and trend analysis.

\subsection{Modular Compliance Engine and GDPR Audit Support} [Estimated: 1,200 words]
A central design principle of this framework is \textbf{modularity with explicit regulatory audit hooks}. The compliance engine is structured as a pluggable module that exposes clear interfaces for GDPR-specific compliance checks, enabling both transparency and extensibility.

The framework treats GDPR compliance not as a monolithic, hard-coded rule set but as a \textbf{configurable audit specification}. This is achieved through:

\begin{enumerate}
\item \textbf{Separation of policy from mechanism}: The \texttt{IComplianceEngine} interface decouples the auditing logic (how data is collected and analysed) from the regulatory rules (what constitutes a violation). The expert baseline library and dynamic threshold parameters are stored in external JSON configuration, allowing GDPR-specific thresholds to be updated without modifying core code.
\item \textbf{Audit trail generation}: Every audit cycle produces a structured \texttt{ComplianceReport} that records: (a) the permission call statistics for each app; (b) the baseline and threshold values applied; (c) the violation status and reason. This report serves as an auditable record that can be presented to data protection authorities, fulfilling GDPR Article 5(2) accountability requirements.
\item \textbf{Regulatory framework parameterisation}: The engine supports dynamic selection of regulatory frameworks (GDPR, PIPL, CCPA) through the \texttt{setRegulatoryFramework()} method. Each framework defines its own mapping of ``special categories'' and permissible thresholds. For GDPR, the engine specifically flags permissions that access health-related data (location, heart rate sensors, activity recognition) as high-sensitivity, applying stricter baseline multipliers.
\item \textbf{Verifiable consent revocation tracking}: The framework monitors not only call frequencies but also the correlation between UI-level permission toggles and actual kernel-level API usage. By comparing the system's recorded permission state with the app's historical call logs, it can detect discrepancies——a direct audit of GDPR Article 7.3 compliance (ease of withdrawal). If an app continues to access a permission after the user has revoked it via system settings, the framework flags a critical violation.
\end{enumerate}

The data access auditing capability builds upon Android 11's (API level 30) data access auditing APIs. These APIs allow applications to register an \texttt{AppOpsManager.OnOpNotedCallback} instance, which performs actions each time private data is accessed. This mechanism enables the framework to gain insights into how apps and their dependencies access private data, helping to identify potentially unexpected data access patterns. By leveraging this callback infrastructure, our compliance engine can perform real-time monitoring of permission usage events without requiring root privileges or VPN-based interception.

\begin{lstlisting}[caption=IComplianceEngine core interface with regulatory audit support]
public interface IComplianceEngine {
    // Load baseline configuration (GDPR/PIPL/CCPA profiles)
    void loadBaseline(Context context, RegulatoryFramework framework);
    
    // Perform audit and produce verifiable report
    ComplianceReport audit(List<HistoricalOp> ops);
    
    // Get current thresholds for verification
    ThresholdConfig getThresholdConfig();
    
    // Switch regulatory framework dynamically
    void setRegulatoryFramework(RegulatoryFramework framework);
    
    // Verify consent revocation effectiveness
    RevocationVerification verifyRevocation(String packageName, String permission);
}
\end{lstlisting}

\subsection{Automated Test Harness Design} [Estimated: 1,000 words]
To systematically validate the framework's detection accuracy, we introduce a dedicated \textbf{Violation Simulator} module. This test harness operates in a controlled environment, generating synthetic permission call logs that mimic both benign and malicious behaviour.

The simulator's architecture comprises three core components:
\begin{enumerate}
\item \textbf{Random Configuration Generator}: Produces test parameters including target permission type, total call count (ranging from 0 to 300), distribution pattern (uniform, bursty, periodic), and time window. This stochastic approach ensures coverage of the legal input space, providing statistical confidence in the engine's performance metrics.
\item \textbf{Scheduler}: Uses WorkManager's \texttt{OneTimeWorkRequest} to dispatch simulated permission calls at the configured random times. OneTimeWorkRequest is designed for scheduling non-repeating work, making it suitable for precisely controlled simulation events. The WorkRequest object encapsulates all information needed for scheduling and running work, including constraints and scheduling information. This ensures the simulation respects Android's background execution limits while providing fine-grained control over event timing.
\item \textbf{Ground Truth Recorder}: Logs the exact parameters of each simulated run, serving as the ``expected result'' against which the audit engine's output is compared during validation.
\end{enumerate}

This design enables N-round independent test cycles (N=50 in our experimental configuration), each with randomised parameters, providing statistical confidence in the engine's precision and recall metrics. The simulator runs entirely within the application's test environment and does not interact with real device sensors or external services.

The stochastic testing methodology is grounded in the principle that algorithmic correctness is a function of the input-output mapping, independent of the physical source of inputs. As long as random samples cover the legal input space—including edge cases such as zero-call scenarios, extremely high-frequency bursts, and permission state transitions—performance on random samples serves as a statistically unbiased estimator of real-world performance. This approach enables rapid iteration and validation without the logistical overhead of deploying to physical devices for every test cycle.

\begin{figure}[H]
\centering
\begin{tikzpicture}[node distance=1.2cm, auto, every node/.style={align=center}]
    % Nodes
    \node (app) [draw, rectangle, fill=bluebox, minimum width=3cm, minimum height=1cm] {Android App \\ (Permission Requests)};
    \node (appops) [draw, rectangle, fill=bluebox, below of=app, node distance=2cm] {AppOpsManager \\ (Historical Ops)};
    \node (work) [draw, rectangle, fill=greenbox, below of=appops, node distance=2cm] {WorkManager \\ (24h Periodic)};
    \node (room) [draw, rectangle, fill=greenbox, right of=work, node distance=4cm] {Room Database \\ (Deduplication)};
    \node (engine) [draw, rectangle, fill=orangebox, below of=work, node distance=2.5cm, minimum width=5cm] {Compliance Engine \\ (Modular + GDPR Audit)};
    \node (sim) [draw, rectangle, fill=redbox, left of=engine, node distance=4cm] {Violation Simulator \\ (Test Harness)};
    \node (monte) [draw, rectangle, fill=redbox, below left of=engine, node distance=2.5cm, xshift=-2.5cm] {Monte Carlo \\ Logical Injection};
    \node (smoke) [draw, rectangle, fill=redbox, below right of=engine, node distance=2.5cm, xshift=2.5cm] {Single Smoke Test \\ (24h real device)};
    \node (report) [draw, rectangle, fill=greenbox, below of=engine, node distance=2.5cm] {Compliance Report \\ (Auditable Log)};

    % Arrows
    \draw[->] (app) -- (appops);
    \draw[->] (appops) -- (work);
    \draw[->] (work) -- (room);
    \draw[->] (room) |- (engine);
    \draw[->] (work) -- (engine);
    \draw[->] (sim) -- (engine);
    \draw[->] (monte) -- (engine);
    \draw[->] (smoke) -- (engine);
    \draw[->] (engine) -- (report);
\end{tikzpicture}
\caption{System architecture overview showing modular compliance engine and test harness}
\label{fig:arch}
\end{figure}

\subsection{Layered Validation Strategy} [Estimated: 500 words]
The framework employs a two-layer validation strategy to ensure algorithmic correctness while maintaining engineering feasibility:

\textbf{Logical layer (Monte Carlo injection)}: bypasses physical hardware, generates large random permission call samples in the JUnit memory environment (random package names, random permission types, random counts 0--300, random timestamps). This completes 1000 rounds of random testing within 3 minutes, validating the algorithmic correctness——including precision, recall, and F1-Score. The theoretical foundation is that algorithmic correctness is a function of the input-output mapping, independent of the physical source of inputs. As long as random samples cover the legal input space, performance on random samples is a statistically unbiased estimator of real-world performance.

\textbf{Integration layer (automated simulation + single smoke test)}: The violation simulator runs N independent 24-hour test cycles, generating randomised call patterns. The audit engine's outputs are automatically compared against ground truth. Additionally, one end-to-end test on a real device uses a specially crafted malicious APK to verify the end-to-end flow and obtain screenshots for the thesis defence.

The integration layer testing specifically addresses the verification of WorkManager's scheduling reliability under real device conditions. PeriodicWorkRequest has a minimum interval of 15 minutes and is guaranteed to run exactly once during each interval, subject to OS battery optimisations such as Doze mode[reference:15]. Our 24-hour test cycles validate that the audit worker executes as scheduled across diverse device states—including Doze mode transitions, background execution restrictions, and varying battery levels—ensuring the framework's reliability in production deployment scenarios.